\section*{Bab \ref{ch:g7:triangles} --- Segitiga dan Sudut}

\begin{solution}{exo:g7:triangles:1}
Penyikunya: $60^\circ$; $45^\circ$; $18^\circ$.
Pelurusnya: $150^\circ$; $135^\circ$; $108^\circ$.
\end{solution}

\begin{solution}{exo:g7:triangles:2}
Sudut yang \omterm{def:g7:triangles:pairs}{bertolak belakang} dengannya juga berukuran $35^\circ$
(\cref{thm:g7:triangles:equalities}); kedua sudut yang tersisa
berpelurus dengannya: $180 - 35 = 145^\circ$ masing-masing.
\end{solution}

\begin{solution}{exo:g7:triangles:3}
(a) $180 - 100 = 80^\circ$. \quad
(b) $180 - 118 = 62^\circ$. \quad
(c) $180 - 150 = 30^\circ$.
\end{solution}

\begin{solution}{exo:g7:triangles:4}
Puncak $40^\circ$: kedua sudut alasnya berbagi $180 - 40 = 140^\circ$
sama rata: $70^\circ$ masing-masing.

Sudut alas $70^\circ$: kedua sudut alasnya sama besar, jadi puncaknya
$180 - 2 \times 70 = 40^\circ$.
\end{solution}

\begin{solution}{exo:g7:triangles:5}
$(6, 8, 12)$: $12 < 6 + 8 = 14$: segitiganya ada.

$(5, 5, 11)$: $11 > 5 + 5 = 10$: mustahil.

$(4, 9, 13)$: $13 = 4 + 9$: ``segitiga'' yang datar --- titiknya
segaris, jadi bukan segitiga sungguhan.

$(7, 7, 7)$: $7 < 14$: \omterm{def:g6:shapes:triangles}{segitiga sama sisi}.
\end{solution}

\begin{solution}{exo:g7:triangles:6}
$x + 2x + 3x = 6x = 180^\circ$, jadi $x = 30^\circ$: sudutnya
$30^\circ$, $60^\circ$ dan $90^\circ$ --- \omterm{def:g6:shapes:triangles}{segitiga siku-siku}.
\end{solution}

\begin{solution}{exo:g7:triangles:7}
Ramalan: $\widehat{BAC} = 180 - (50 + 60) = 70^\circ$; busur derajat
pada gambar yang sudah jadi membenarkannya.
\end{solution}

\begin{solution}{exo:g7:triangles:8}
Pada perpotongan yang pertama, keempat sudutnya adalah $54^\circ$,
$126^\circ$, $54^\circ$, $126^\circ$ (pasangan \omterm{def:g7:triangles:pairs}{bertolak belakang} dan
pelurusnya). \omterm{def:g6:lines:perp}{Garis sejajarnya} mengulang keempat ukuran yang sama pada
perpotongan yang kedua (sudut sehadap): seluruhnya delapan sudut,
empat berukuran $54^\circ$ dan empat berukuran $126^\circ$.
\end{solution}

\begin{solution}{exo:g7:triangles:9}
Misalkan $c$ sisi ketiganya. Ketaksamaan segitiga menuntut
$c < 8 + 3 = 11$ dan juga $8 < c + 3$, yaitu $c > 5$. Jadi
$5 < c < 11$: panjang $7$ cm berhasil; panjang $12$ cm (atau $4$ cm)
tidak.
\end{solution}

\begin{solution}{exo:g7:triangles:10}
$\widehat A + \widehat B = 180 - 90 = 90^\circ$, dan
$\widehat A = 2\widehat B$, jadi $3\widehat B = 90^\circ$:
$\widehat B = 30^\circ$ dan $\widehat A = 60^\circ$.
\end{solution}

\begin{solution}{exo:g7:triangles:11}
Diagonal $[AC]$ membelah $ABCD$ menjadi segitiga $ABC$ dan $ACD$.
Keempat sudut \omterm{def:g4:geometry:polygon}{segi empatnya} persis keenam sudut kedua segitiga itu,
yang dikelompokkan ulang (sudut di $A$ dan di $C$ masing-masing
terbelah dua). \omterm{def:g1:addition:def}{Jumlahnya}: $180 + 180 = 360^\circ$. \omterm{def:g4:geometry:polygon}{Segi lima} terbelah
dari satu \omterm{def:g5:solids:def}{titik sudut} menjadi tiga segitiga: $3 \times 180 = 540^\circ$.
\end{solution}

\begin{solution}{pb:g7:triangles:1}
\textbf{1.} Diagonal dari satu \omterm{def:g5:solids:def}{titik sudut} \omterm{def:g4:geometry:polygon}{segi enam} memotongnya
menjadi $4$ segitiga, jadi sudut dalamnya berjumlah
$4 \times 180 = 720^\circ$.

\textbf{2.} Dari satu \omterm{def:g5:solids:def}{titik sudut}, diagonalnya menuju semua \omterm{def:g5:solids:def}{titik
sudut} kecuali dirinya sendiri dan kedua tetangganya: $n - 3$ diagonal,
yang memotong \omterm{def:g4:geometry:polygon}{segi banyaknya} menjadi $n - 2$ segitiga. Setiap sudut
dalam \omterm{def:g4:geometry:polygon}{segi banyaknya} terbagi di antara segitiganya, jadi \omterm{def:g1:addition:def}{jumlahnya}
\[
(n - 2) \times 180^\circ .
\]
Periksa: $n = 3$: $180^\circ$; $n = 4$: $360^\circ$
(\cref{exo:g7:triangles:11}); $n = 5$: $540^\circ$; $n = 6$:
$720^\circ$.

\textbf{3.} Segi sepuluh: $8 \times 180 = 1\,440^\circ$. Segi dua
belas: $10 \times 180 = 1\,800^\circ$.

\textbf{4.} Dengan membagi $n$: segitiga $60^\circ$; persegi
$90^\circ$; \omterm{def:g4:geometry:polygon}{segi lima} $\frac{540}{5} = 108^\circ$; \omterm{def:g4:geometry:polygon}{segi enam}
$\frac{720}{6} = 120^\circ$; \omterm{def:g4:geometry:polygon}{segi delapan} $\frac{1080}{8} =
135^\circ$.

\textbf{5.} Menambah satu \omterm{def:g5:solids:def}{titik sudut} menambah satu segitiga lagi pada
kipasnya: \omterm{def:g4:geometry:polygon}{segi banyak} yang baru perlu $n - 1$ segitiga sedangkan yang
lama perlu $n - 2$. Satu segitiga tambahan, satu $180^\circ$ tambahan.

\textbf{6.} Pejalannya berbelok sejauh pelurusnya:
$180^\circ - \widehat A$. Untuk segitiga $40$--$60$--$80$, belokannya
$140^\circ$, $120^\circ$ dan $100^\circ$.

\textbf{7.} $140 + 120 + 100 = 360^\circ$: satu putaran penuh.

\textbf{8.} Sepanjang seluruh perjalanan, hidungmu bermula dan
berakhir menghadap ke arah yang sama, setelah menyapu tepat satu kali
\omterm{def:g6:measure:perimeter}{keliling}: seluruh belokan yang dilakukan di pojoknya berjumlah satu
putaran lengkap, $360^\circ$ --- berapa pun banyak pojoknya dan apa
pun bentuk bloknya. (Kenyataan itu indahnya tidak bergantung pada
$n$.)

\textbf{9.} Pada setiap pojok, sudut dalam $+$ belokan $= 180^\circ$.
Menjumlahkan atas $n$ pojoknya:
\[
\text{(jumlah sudut dalam)} + 360^\circ = 180^\circ \times n,
\qquad\text{jadi}\qquad
\text{jumlah sudut dalam} = (n - 2) \times 180^\circ
\]
--- rumus pertanyaan 2, yang dibuktikan ulang dengan berjalan.

\textbf{10.} Sudut dalam $150^\circ$ berarti setiap belokannya
$30^\circ$, dan $n = \frac{360}{30} = 12$: segi dua belas beraturan.
Sudut dalam $155^\circ$ akan berarti belokan $25^\circ$, dan
$\frac{360}{25} = 14.4$ bukan banyaknya pojok yang bulat: \omterm{def:g4:geometry:polygon}{segi banyak}
beraturan seperti itu tidak ada.

\textbf{11.} Segitiga ($60^\circ$): $6$ bertemu,
$6 \times 60 = 360$. Persegi ($90^\circ$): $4$ bertemu. \omterm{def:g4:geometry:polygon}{Segi enam}
($120^\circ$): $3$ bertemu, $3 \times 120 = 360$. Ketiganya menutup
dengan sempurna.

\textbf{12.} Tiga pojok \omterm{def:g4:geometry:polygon}{segi lima}: $3 \times 108 = 324^\circ$ ---
tersisa celah $36^\circ$. Empat pojoknya:
$4 \times 108 = 432^\circ > 360^\circ$ --- tumpang tindih. Tidak ada
yang berhasil: \omterm{def:g4:geometry:polygon}{segi lima} beraturan tidak dapat mengubin bidang.

\textbf{13.} Pada satu pojok pengubinan sedikitnya $3$ ubin bertemu.
\omterm{def:g4:geometry:polygon}{Segi banyak} beraturan dengan lebih dari enam sisi punya sudut dalam
yang \emph{lebih besar} daripada $120^\circ$ (pertanyaan 4 dan
tumbuhnya sudut itu bersama $n$), jadi tiga pojoknya melampaui
$3 \times 120 = 360^\circ$: tumpang tindih, mustahil. Dengan enam sisi
tepat, $3 \times 120 = 360$ berhasil; di bawah enam, pertanyaan 11
dan 12 memilah keadaannya. Daftar lengkap pengubin beraturannya:
segitiga, persegi, \omterm{def:g4:geometry:polygon}{segi enam}.

\textbf{14.} \omterm{def:g4:geometry:polygon}{Segi delapan} dan persegi:
$135 + 135 + 90 = 360^\circ$: pojoknya menutup --- pola jalanan yang
klasik. Persegi, \omterm{def:g4:geometry:polygon}{segi enam}, segi dua belas: belokan segi dua belasnya
$\frac{360}{12} = 30^\circ$, jadi sudut dalamnya $150^\circ$, dan
$90 + 120 + 150 = 360^\circ$: itu berhasil juga.

\textbf{15.} Pojok sarang lebah: $3 \times 120 = 360^\circ$,
sempurna. Di antara ubin segitiga, persegi dan \omterm{def:g4:geometry:polygon}{segi enam} yang sama
luasnya, \omterm{def:g4:geometry:polygon}{segi enamnya} punya \omterm{def:g6:measure:perimeter}{keliling} terpendek --- ia yang paling
dekat di antara ketiganya dengan lingkaran Dido
(\cref{pb:g6:measure:1}) --- jadi sel \omterm{def:g4:geometry:polygon}{segi enam} menampung madu yang
sama dengan lilin paling sedikit: lebah mengubin seperti ahli geometri.

\end{solution}
